\chapter{Trao đổi yếu có khối lượng, các hạt ảo và đối xứng bị phá vỡ}

Nhóm của Mô hình Chuẩn trở nên bớt bí ẩn đi nhiều khi ta đọc nó như một tập các chuyển động được phép trên bảng fermion. Các phép biến đổi màu dịch chuyển theo phương dọc giữa đỏ, lục và lam. Các phép biến đổi yếu tích điện dịch chuyển theo phương ngang trong các bộ đôi quark và lepton. Các phép biến đổi pha điện từ cung cấp một thao tác thứ ba. Các boson chuẩn tương ứng không phải là những thứ trang trí được thêm vào sau khi các đối xứng đã được tìm ra: việc phát ra một boson chuẩn về mặt vật lý thực hiện đúng cùng một sự thay đổi trạng thái mà một phần tử sinh của đối xứng mô tả.

Nhận xét đó đặt ra câu hỏi chính của bài giảng này. Một đỉnh yếu có hằng số liên kết cỡ tương đương một đỉnh điện từ, thế nhưng phân rã beta lại chậm hơn khổng lồ so với một quá trình mạnh hay điện từ điển hình. Mảnh ghép còn thiếu là sự truyền. Ở xung lượng thấp, việc trao đổi một \(W\) có khối lượng phải trả một thừa số cỡ \(1/m_W^2\) trong biên độ. Trước hết ta sẽ phát biểu ấy trong không gian khoảng cách rồi trong không gian xung lượng, theo các hệ quả của nó qua một cuộc thảo luận dài về các hạt thực và hạt ảo, và cuối cùng chuyển sang mô hình định tính đầu tiên về sự phá vỡ đối xứng tự phát.

\section{Đọc đối xứng chuẩn từ bảng fermion}

\subsection{Màu tác dụng theo phương dọc}

Sắp xếp sáu hương vị quark thành ba họ yếu,
\begin{equation}
  (u,d),\qquad(c,s),\qquad(t,b),
\end{equation}
và gán cho mỗi quark một nhãn màu đỏ, lục hoặc lam. Bảng quark thu được có ba hàng màu và sáu cột hương vị. \(SU(3)\) màu trộn ba phần tử trong bất kỳ một cột hương vị nào. Chẳng hạn, một phần tử sinh có thể biến một quark up đỏ thành một quark up lục. Trong ngôn ngữ hạt, cùng thao tác ấy đi kèm với việc phát ra một gluon mang màu và phản màu bù trừ.

Cùng một phép biến đổi màu phải tác dụng lên mọi hương vị. Gluon biến một quark up đỏ thành một quark up lục cũng là đúng loài gluon biến một quark top đỏ thành một quark top lục. Một phép toán đối xứng đích thực không thể chỉ áp dụng cho thế hệ thứ nhất mà bỏ qua các thế hệ khác. Vì thế chuyển động màu là theo phương dọc và phổ quát:
\begin{equation}
  q_r\longleftrightarrow q_g\longleftrightarrow q_b,
  \qquad q\in\{u,d,c,s,t,b\}.
\end{equation}
Nó làm thay đổi màu nhưng không làm thay đổi hương vị.

Đó là mối liên hệ đặc trưng giữa một nhóm chuẩn và các boson chuẩn của nó. Các phần tử sinh vi phân mô tả các phép biến đổi giữa các trạng thái, còn sự phát ra hoặc hấp thụ các boson tương ứng hiện thực hóa các phép biến đổi ấy thành những quá trình vật lý.

\subsection{Các phép biến đổi yếu tác dụng theo phương ngang}

Tương tác yếu tích điện dịch chuyển theo hướng kia. Một quark kiểu up có thể phát ra một boson yếu tích điện và trở thành đối tác kiểu down của nó mà không đổi màu:
\begin{equation}
  u_b\longrightarrow d_b+W^+.
\end{equation}
Chính cùng \(W^+\) ấy cũng xuất hiện trong
\begin{equation}
  c_g\longrightarrow s_g+W^+,
  \qquad
  t_r\longrightarrow b_r+W^+.
\end{equation}
Các nhãn \(r,g,b\) chỉ là khán giả. Gluon không dịch chuyển theo phương ngang qua bảng, còn các boson \(W\) tích điện không dịch chuyển theo phương dọc qua màu.

Thêm đối xứng pha điện từ \(U(1)\), dưới đó mọi trường mang điện đều nhận một pha. Khi ấy cấu trúc nhóm ở mức đầu tiên là
\begin{equation}
  SU(3)\times SU(2)\times U(1).
\end{equation}
Các dấu nhân chéo có nghĩa là đây là các thừa số độc lập: một thừa số tác dụng lên màu, một thừa số tác dụng lên các bộ đôi yếu, và một thừa số tác dụng bằng một pha. Đây là bức tranh sơ bộ hữu ích, nhưng vẫn chưa phải là toàn bộ câu chuyện điện yếu.\editorialnote{Nhóm chuẩn chính xác của Mô hình Chuẩn theo quy ước được viết là
\(SU(3)_c\times SU(2)_L\times U(1)_Y\). \(SU(2)_L\) yếu chỉ tác dụng lên các bộ đôi tay trái, và \(U(1)_{\rm em}\) điện từ cuối cùng xuất hiện từ \(SU(2)_L\times U(1)_Y\) sau khi đối xứng bị phá vỡ. Bài giảng nhấn mạnh rõ rằng cách đồng nhất sơ bộ hai thừa số cuối này vẫn chưa hoàn toàn nhất quán; các bài giảng sau sẽ bổ sung cấu trúc chiral và điện yếu còn thiếu.}

\subsection{Cùng một khuôn mẫu cũng bao gồm các lepton}

Bên dưới ba hàng màu, đặt một hàng lepton không màu:
\begin{equation}
\begin{array}{c|cc|cc|cc}
 &u&d&c&s&t&b\\ \hline
Q&+\frac23&-\frac13&+\frac23&-\frac13&+\frac23&-\frac13\\
\text{lepton partner}
 &\nu_e&e^-&\nu_\mu&\mu^-&\nu_\tau&\tau^-\\
Q&0&-1&0&-1&0&-1
\end{array}
\end{equation}
Trong mỗi cặp kề nhau, các điện tích khác nhau đúng một đơn vị. Đó là lý do cùng một boson yếu tích điện có thể nối \(u\) với \(d\), \(c\) với \(s\), \(t\) với \(b\), và một lepton tích điện với neutrino của nó. Đỉnh \(W\) quan tâm đến các số lượng tử thích hợp và độ chênh điện tích, chứ không quan tâm fermion ấy là quark hay lepton.

\begin{figure}[htbp]
\centering
\includegraphics[width=0.88\textwidth]{lecture_06_figure_07.png}
\caption{Bảng fermion hoàn chỉnh: ba hàng màu quark ở trên hàng lepton không màu, với các điện tích quark và ba bộ đôi lepton được hiển thị.}
\lectureframe{00:18:55}
\end{figure}

Muon có cùng điện tích và các số lượng tử nội tại đã biết như electron, nhưng nặng hơn khoảng hai trăm lần. Tau còn nặng hơn nữa. Khối lượng neutrino rất nhỏ và trong phần thảo luận của bài giảng chưa được biết chi tiết. Ở đây vẫn chưa xuất hiện một nguyên lý nào giải thích được khuôn mẫu khối lượng quan sát được của các họ.

\begin{classroomqa}
\audiencequestion{Vì sao lepton, baryon và meson lại được gọi bằng những tên ấy?\lecturetimestamp{00:19:02}}
\lecturerresponse{Về mặt lịch sử, các gốc từ Hy Lạp ấy có nghĩa gần đúng là nhẹ, nặng và trung gian. Kiểu đặt tên dựa trên khối lượng này về sau trở nên gây hiểu lầm. Muon từng được gọi là mu meson vì khối lượng của nó nằm giữa khối lượng electron và proton, nhưng các tương tác của nó cho thấy nó là một lepton, giống electron hơn nhiều so với pion. Ngày nay, meson có nghĩa là một hadron quark--phản quark và là một boson; muon không phải là meson.}
\end{classroomqa}

\section{Câu đố về tính yếu}

\subsection{Một quá trình dòng tích điện cụ thể}

Dùng phân rã beta của neutron làm giản đồ làm việc. Các thành phần quark hoá trị là
\begin{equation}
  n=udd,\qquad p=uud.
\end{equation}
Một quark down thực hiện chuyển tiếp
\begin{equation}
  d\longrightarrow u+W^-,
\end{equation}
và boson tích điện tạo ra
\begin{equation}
  W^-\longrightarrow e^-+\bar\nu_e.
\end{equation}
Gộp lại,
\begin{equation}
  n\longrightarrow p+e^-+\bar\nu_e.
\end{equation}
Hai quark hoá trị còn lại đi xuyên suốt như các khán giả. Ở đỉnh thứ nhất, neutron trở thành proton; ở đỉnh thứ hai, chính tương tác yếu nối một bộ đôi lepton đã biến \(W^-\) thành electron và phản neutrino electron quan sát được.

\subsection{Ba hằng số liên kết, ba ấn tượng rất khác nhau}

Trong QED, điện tích electron \(e\) là biên độ của đỉnh. Bình phương nó cho ta xác suất phát xạ đặc trưng, với các thừa số quy ước được gom vào
\begin{equation}
  \alpha=\frac{e^2}{4\pi}\simeq\frac1{137}
\end{equation}
trong đơn vị tự nhiên. Tương tác mạnh có hằng số liên kết \(g_s\), và ở một thang hadron điển hình, đại lượng tương tự \(\alpha_s\) của nó có thể cỡ \(0.2\).
Giá trị đó lớn hơn nhiều so với \(1/137\), nên việc phát gluon là chuyện thường gặp khi một quark bị tác động.

Dùng \(g_w\) để ký hiệu hằng số liên kết yếu tích điện. Bình phương của nó không hề nhỏ một cách phi thường; ở mức bậc độ lớn được dùng ở đây, nó tương đương với hằng số liên kết điện từ. Thế nhưng một neutron tự do phải chờ xấp xỉ mười lăm phút mới phân rã, trong khi một phân rã mạnh có thể xảy ra trên thang thời gian băng qua hadron. Ước lượng lịch sử trong bài giảng dùng con số xấp xỉ cũ là mười hai phút, nhưng dù dùng số nào thì câu đố ấy vẫn không thể tránh khỏi.\editorialnote{Thời gian sống trung bình hiện đại của neutron tự do vào khoảng
\(8.8\times10^2\,{\rm s}\), với sự căng thẳng thực nghiệm vẫn đang tiếp diễn giữa các phương pháp đo chính. Con số thời gian sống xấp xỉ trong bài giảng được giữ lại như một thang sư phạm, chứ không phải như một giá trị chính xác.}
\begin{quote}
Nếu \(g_w\) không nhỏ một cách đặc biệt, thì điều gì làm cho các quá trình yếu trở nên yếu?
\end{quote}

Câu trả lời không thể chỉ nằm ở riêng các đỉnh. Mọi giản đồ Feynman còn chứa cả các hàm truyền, tức những thừa số mô tả sự truyền trung gian từ đỉnh này sang đỉnh khác.

\section{Hàm truyền có khối lượng trong không gian khoảng cách}

\subsection{Các đỉnh không phải là toàn bộ giản đồ}

Giản đồ phân rã beta chứa hai đỉnh yếu, đóng góp hai lũy thừa của \(g_w\), và một đường \(W\) nối chúng. Ký hiệu hàm truyền trong không gian vị trí là \(\Delta(x-y)\), trong đó \(x\) là sự kiện phát xạ và \(y\) là sự kiện hấp thụ hoặc chuyển hóa. Nó là hạt nhân biên độ gắn với việc mang nhiễu động trường giữa các sự kiện ấy.

Hàm truyền tương đối tính chính xác chứa nhiều cấu trúc hơn mức ta cần. Đặc trưng định tính quan trọng của nó là sự phụ thuộc vào khối lượng. Ở khoảng cách kiểu không gian ngắn, nó có một kỳ dị dạng lũy thừa, sơ đồ là
\begin{equation}
  \Delta(d)\sim\frac1{d^2}.
\end{equation}
Một trường không khối lượng vẫn giữ một đuôi dạng lũy thừa. Ngược lại, một trường có khối lượng bị dập tắt theo hàm mũ một khi độ tách trở nên lớn so với thang Compton của nó. Vì vậy đồ thị ngoặt xuống rất mạnh.

\subsection{Đơn vị tự nhiên xác định vị trí chỗ ngoặt}

Đặt
\begin{equation}
  c=\hbar=1.
\end{equation}
Khi đó
\begin{equation}
  [m]=[E]=[p]=L^{-1}.
\end{equation}
Các hệ thức quen thuộc giải thích từng dấu bằng: \(E=mc^2\) đồng nhất khối lượng với năng lượng, còn \(p\,x\) có thứ nguyên của \(\hbar\), nên xung lượng và khoảng cách có đơn vị nghịch đảo nhau. Vì thế độ dài duy nhất có thể tạo thành từ khối lượng \(m\) là
\begin{equation}
  d_m\sim\frac1m.
\end{equation}
Khôi phục các hằng số cho ta
\begin{equation}
  d_m\sim\frac{\hbar}{mc},
\end{equation}
chính là bước sóng Compton. Hạt trao đổi càng nặng thì tầm tác dụng hiệu dụng càng ngắn.

\begin{figure}[htbp]
\centering
\includegraphics[width=0.78\textwidth]{lecture_06_figure_02.png}
\caption{Các thứ nguyên trong đơn vị tự nhiên và độ suy giảm định tính của hàm truyền. Đường cong đen tiếp tục với một đuôi lũy thừa; đường cong đỏ có khối lượng ngoặt xuống gần thang nghịch đảo khối lượng.}
\lectureframe{00:27:19}
\end{figure}

\begin{figure}[htbp]
\centering
\begin{tikzpicture}[x=1.0cm,y=0.9cm]
  \draw[->] (0,0) -- (7,0) node[right] {$d$};
  \draw[->] (0,0) -- (0,5) node[above] {$|\Delta(d)|$};
  \draw[thick]
    plot[smooth] coordinates {(0.45,4.5) (0.9,3.3) (1.6,2.5)
      (2.5,1.9) (3.7,1.45) (5.1,1.12) (6.4,0.9)};
  \draw[thick,red]
    plot[smooth] coordinates {(0.45,4.5) (0.9,3.3) (1.6,2.5)
      (2.5,1.9) (3.2,1.55) (3.7,0.95) (4.2,0.36) (4.8,0.1)};
  \draw[dashed] (3.4,0) -- (3.4,1.4);
  \node at (5.7,1.4) {$m=0$};
  \node[red] at (4.75,0.72) {$m>0$};
  \node[below] at (3.4,0) {$d\sim m^{-1}$};
\end{tikzpicture}
\caption{Một sơ đồ minh họa gọn cho lập luận trên bảng. Đây không phải là công thức chính xác cho một hàm truyền bốn chiều; nó tách riêng tầm tác dụng do khối lượng chi phối.}
\end{figure}

\begin{classroomqa}
\audiencequestion{Trục tung biểu diễn đại lượng gì?
\lecturetimestamp{00:28:30}}
\lecturerresponse{Độ lớn của hàm truyền, ký hiệu là \(\Delta\). Bình phương của nó đóng góp vào một xác suất sau khi đã bao gồm biên độ đầy đủ, các trạng thái ngoài và các thừa số không gian pha. Đồ thị nhằm cho thấy biên độ của sự trao đổi có khối lượng trên một khoảng cách lớn tắt đi rất nhanh.\editorialnote{Bản thân một hàm truyền không phải là một biên độ xác suất đo trực tiếp được cho một hạt đi theo một quỹ đạo cổ điển. Cách nói trong bài giảng là một ước lượng hữu ích về tầm tác dụng: hàm truyền là một thừa số trong biên độ tán xạ hoặc phân rã đầy đủ.}}
\end{classroomqa}

\section{Biến đổi Fourier và giới hạn Fermi}

\subsection{Thí nghiệm xác định các xung lượng}

Một thí nghiệm thường đo các xung lượng ngoài hơn là hai điểm không-thời gian nơi một đường trong bắt đầu và kết thúc. Trong phân rã beta, xung lượng do \(W\) mang được suy ra từ hiệu giữa xung lượng của neutron và proton, và nó bằng tổng xung lượng của electron và phản neutrino. Hàm truyền trong không gian vị trí và trong không gian xung lượng liên hệ với nhau qua một biến đổi Fourier.
\begin{equation}
  k=p_n-p_p=p_e+p_{\bar\nu_e}.
\end{equation}

Trong quy ước mẫu số dương được dùng trên bảng,
\begin{equation}
  \widetilde{\Delta}_W(k)=\frac1{k^2+m_W^2},
  \qquad
  \widetilde{\Delta}_\gamma(k)=\frac1{k^2}.
\end{equation}
Hàm truyền photon không khối lượng trở nên lớn khi độ truyền xung lượng nhỏ. Hàm truyền \(W\) có khối lượng không bao giờ vượt quá một thang cỡ \(1/m_W^2\). Đối với phân rã beta,
\begin{equation}
  |k^2|\ll m_W^2,
\end{equation}
nên các chi tiết của \(k^2\) là không đáng kể và
\begin{equation}
  \widetilde{\Delta}_W(k)\simeq\frac1{m_W^2}.
\end{equation}

\begin{classroomqa}
\audiencequestion{Trong \(k^2\) tương đối tính có nên có một dấu trừ giữa năng lượng và xung lượng không gian không?
\lecturetimestamp{00:35:11}}
\lecturerresponse{Có. Với metric Minkowski chủ yếu dấu trừ,
\(k^2=(k^0)^2-\boldsymbol{k}^2\), và mẫu số vô hướng theo quy ước là
\(k^2-m_W^2+i\epsilon\). Trên bảng đang dùng độ lớn kiểu không gian (spacelike) dương, hay tương đương là độ lớn Euclid dương. Trong phân rã beta năng lượng thấp, bốn-xung lượng truyền đi là rất nhỏ so với \(m_W\), nên dù dùng ký hiệu nào thì cũng cho cùng sự dập tắt \(1/m_W^2\), sai khác chỉ bởi một dấu không quan trọng trong biên độ.}
\end{classroomqa}

\subsection{Biên độ năng lượng thấp}

Hai đỉnh yếu và một hàm truyền cho
\begin{align}
  \mathcal A_{\rm weak}
  &\sim g_w\,\widetilde{\Delta}_W(k)\,g_w\\
  &\sim\frac{g_w^2}{k^2+m_W^2}
  \simeq\frac{g_w^2}{m_W^2}.
\end{align}
Bình phương biên độ tạo ra sự dập tắt đặc trưng
\begin{equation}
  |\mathcal A_{\rm weak}|^2
  \sim\frac{g_w^4}{m_W^4},
\end{equation}
trước khi đưa vào các thừa số động học và không gian pha thích hợp cho phân rã cụ thể đang xét. Tương tác yếu là yếu ở các năng lượng hạt nhân không phải vì hằng số liên kết chuẩn cơ bản của nó nhỏ đến mức phi lý, mà vì hạt trung gian của nó nặng.

\begin{figure}[htbp]
\centering
\includegraphics[width=0.86\textwidth]{lecture_06_figure_03.png}
\caption{Quá trình yếu neutron--proton và độ truyền xung lượng do đường trong mang, như đã được thảo luận khi xác định dấu của hàm truyền.}
\lectureframe{00:35:40}
\end{figure}

\begin{figure}[htbp]
\centering
\includegraphics[width=0.84\textwidth]{lecture_06_figure_04.png}
\caption{Cấu trúc phân rã beta bên cạnh ước lượng năng lượng thấp được đóng khung \(g_w^2/m_W^2\).}
\lectureframe{00:37:25}
\end{figure}

Về mặt số trị, \(m_W\) xấp xỉ 100 lần khối lượng proton ở mức chính xác của bài giảng này. Một thừa số \(m_W^{-2}\) trong biên độ trở thành \(m_W^{-4}\) trong tốc độ, nên ngay cả một hằng số ghép có cỡ thông thường của lý thuyết gauge cũng tạo ra một thời gian sống rất dài.

\begin{classroomqa}
\audiencequestion{Ta đang nói đến độ suy giảm cỡ một phần triệu hay một phần tỷ, và ban đầu người ta biết khối lượng của \(W\) bằng cách nào?
\lecturetimestamp{00:38:11}}
\lecturerresponse{Phép đếm thô của bài giảng gợi ra rằng xác suất yếu ở năng lượng thấp nhỏ ở mức cỡ một phần tỷ so với một quá trình hadron không bị suy giảm. Về mặt lịch sử, các thí nghiệm năng lượng thấp không đo riêng rẽ \(g_w\) và \(m_W\). Chúng đo tổ hợp nhỏ \(g_w^2/m_W^2\), được gọi là hằng số Fermi. Một khi hằng số ghép gauge được gắn với lý thuyết điện yếu, tổ hợp đó dự đoán một hạt trung gian rất nặng; về sau các máy gia tốc đã trực tiếp tạo ra \(W\) và đo khối lượng của nó.}
\end{classroomqa}

Ở các xung lượng nhỏ hơn rất nhiều so với \(m_W\), mọi thí nghiệm đều thấy gần như cùng một hằng số:
\begin{equation}
  G_F\sim\frac{g_w^2}{m_W^2}.
\end{equation}
Đó là vì sao Fermi có thể mô tả phân rã beta trong những năm 1930 như một tương tác điểm kiểu bốn-fermion, rất lâu trước khi hạt trung gian được quan sát. Khi năng lượng va chạm tăng lên, các thí nghiệm trở nên nhạy với sự phụ thuộc theo xung lượng ẩn trong hàm truyền, qua đó xác nhận rằng hằng số hiệu dụng ấy là cái bóng năng lượng thấp của sự trao đổi hạt nặng.\editorialnote{Với quy ước chuẩn hóa điện yếu tiêu chuẩn ở mức cây,
\[
  \frac{G_F}{\sqrt2}=\frac{g^2}{8m_W^2}.
\]
Bài giảng cố ý lược bỏ các hệ số số học và spinor, và chỉ giữ lại luật tỉ lệ \(G_F\sim g_w^2/m_W^2\). \(W\) và \(Z\) đã được phát hiện trực tiếp tại CERN vào 1983; mốc niên đại thoáng qua vào khoảng giữa đến cuối thập niên 1970 của bài giảng là quá sớm đối với quan sát trực tiếp.}

Về mặt lịch sử, ý tưởng rằng các tương tác yếu có thể được trung gian bởi một hạt nặng xuất hiện vào giữa thế kỷ 20; bài giảng nhắc đến Julian Schwinger trong số các nhà lý thuyết đã khảo sát khả năng đó và lưu ý rằng chữ cái \(W\) có lẽ là viết tắt của weak. Dữ liệu năng lượng thấp đã ngụ ý rằng một hạt trung gian như vậy hẳn phải cực kỳ nặng trên thang hạt nhân, từ rất lâu trước khi các máy gia tốc có đủ năng lượng để tạo ra nó như một hạt thực.

\subsection{Ba yếu tố độc lập chi phối một tốc độ}

Một câu hỏi trên lớp về liên kết hạt nhân tách ra ba yếu tố điều khiển một chuyển tiếp, về mặt logic là khác nhau:
\begin{enumerate}
  \item các hằng số ghép tại đỉnh;
  \item các hàm truyền, bao gồm khối lượng và sự phụ thuộc theo xung lượng;
  \item năng lượng sẵn có từ trạng thái đầu tới trạng thái cuối và không gian pha.
\end{enumerate}
Liên kết có thể đóng một kênh phân rã vì trạng thái cuối sẽ có năng lượng lớn hơn trạng thái đầu. Đó là một sự cấm đoán động học, không phải là sự giảm thêm của \(g_w\) hay sự thay đổi của \(m_W\). Khi một kênh vẫn còn mở nhưng có rất ít năng lượng sẵn có, không gian pha của nó vẫn có thể bị suy giảm mạnh.

\begin{classroomqa}
\audiencequestion{Trường hay liên kết của một proton ở gần có làm suy giảm phân rã beta thông qua cùng cơ chế hàm truyền không?
\lecturetimestamp{00:41:13}}
\lecturerresponse{Không phải trong ví dụ đang được bàn tới. Liên kết quan trọng vì nó làm thay đổi năng lượng của các trạng thái hạt nhân đầu và cuối. Một quá trình không thể xảy ra nếu trạng thái cuối không thể tiếp cận được về mặt năng lượng. Hằng số ghép, hàm truyền và năng lượng sẵn có là ba thành phần tách biệt.}
\end{classroomqa}

\begin{classroomqa}
\audiencequestion{Các hằng số ghép có thật sự là hằng số không?
\lecturetimestamp{00:42:55}}
\lecturerresponse{Chỉ ở một thang cố định. Các hiệu chỉnh lượng tử làm cho các hằng số ghép hiệu dụng chạy theo xung lượng hoặc khoảng cách. Lập luận hiện tại xét mỗi hằng số ghép ở thang thích hợp với quá trình; sự phụ thuộc theo thang sẽ được phát triển sau.}
\end{classroomqa}

\section{Trao đổi ảo mà không từ bỏ các định luật bảo toàn}

\subsection{Vì sao đường của hạt nặng có thể xuất hiện}

Phân rã neutron thông thường không có đủ năng lượng để tạo ra một \(W\) tự do, ở trên vỏ khối. \(W\) trong giản đồ là một đường nội và ở ngoài vỏ khối. Bài giảng giới thiệu sự kiện này bằng ước lượng năng lượng--thời gian truyền thống
\begin{equation}
  \Delta E\,\Delta t\lesssim\hbar:
\end{equation}
một cấu hình trung gian mà năng lượng cổ điển của nó dường như lệch đi một lượng \(\Delta E\) chỉ có thể đóng góp trong một khoảng thời gian cỡ \(\hbar/\Delta E\).

Phép loại suy đường hầm làm cho hình dung trực quan trở nên sống động. Đặt một hạt alpha trong một giếng thế hạt nhân. Năng lượng của nó nằm dưới đỉnh rào, nên theo cơ học cổ điển nó không thể thoát ra. Về mặt lượng tử, hàm sóng kéo dài xuyên qua rào, và hạt có thể xuất hiện tại một điểm ở phía bên kia với cùng tổng năng lượng như ban đầu. Nếu ta cố mô tả phần bên trong rào như một quỹ đạo cổ điển thông thường, thì thế năng dường như quá lớn và động năng sẽ buộc phải âm. Bức tranh cổ điển thất bại đó tương tự với việc cố coi \(W\) bên trong là một hạt thực.

\begin{figure}[htbp]
\centering
\includegraphics[width=0.78\textwidth]{lecture_06_figure_06.png}
\caption{Phép loại suy đường hầm: điểm ban đầu và điểm được đánh dấu trên sườn phía xa có cùng tổng năng lượng, còn rào thế xen giữa bị cấm theo cổ điển.}
\lectureframe{00:46:00}
\end{figure}

\begin{figure}[htbp]
\centering
\begin{tikzpicture}[x=1.0cm,y=0.9cm]
  \draw[->] (0,0) -- (7,0) node[right] {$x$};
  \draw[->] (0,0) -- (0,4.8) node[above] {$V(x)$};
  \draw[thick]
    plot[smooth] coordinates {(0.3,1.4) (1.0,0.7) (1.7,0.55)
      (2.1,0.9) (3.2,3.9) (4.2,2.1) (5.0,0.75) (6.6,0.18)};
  \draw[dashed] (1.7,0.75) -- (5.0,0.75);
  \fill (1.7,0.75) circle (0.06);
  \fill (5.0,0.75) circle (0.06);
  \node at (3.35,0.35) {$\Delta E\,\Delta t\lesssim\hbar$};
\end{tikzpicture}
\caption{Một bản dựng lại dễ đọc của rào thế và các điểm có cùng năng lượng.}
\end{figure}

Trong bức tranh phân rã beta, neutron có thể thăng giáng thành một proton cộng với một \(W^-\) nội, nhưng đường nặng ấy phải nhanh chóng được khử đi. Nó có thể bị tái hấp thụ, khôi phục neutron, hoặc có thể kết thúc ở đỉnh thứ hai và tạo ra electron cùng phản neutrino nhẹ hơn rất nhiều. Tổng năng lượng của trạng thái đầu và cuối là như nhau.\editorialnote{Năng lượng và xung lượng được bảo toàn tại mọi đỉnh Feynman.
Các hạt ảo không thật sự vay năng lượng rồi trả lại về sau. Phát biểu chính xác là một đường trong được lấy tích phân trên các bốn-xung lượng không nhất thiết phải thỏa mãn hệ thức vỏ khối của hạt thực
\(k^2=m_W^2\). Ngôn ngữ nguyên lý bất định và đường hầm được giữ lại vì đó là lộ trình giải thích của bài giảng, nhưng phát biểu ngoài vỏ khối mới là cách diễn đạt hiện đại của QFT.}

\subsection{Các boson W thực}

Một máy gia tốc đủ năng lượng có thể tạo ra một \(W\) thực. Khi đó nó có thời gian sống riêng của chính nó, không liên quan tới một thời gian bất định giả định nào của đường ảo. Một \(W\) thực đang chuyển động sẽ sống lâu hơn trong phòng thí nghiệm do giãn thời gian tương đối tính thông thường. Trong phân rã lepton hai vật của một \(W\) đứng yên,
\begin{equation}
  W^-\longrightarrow e^-+\bar\nu_e,
\end{equation}
electron và phản neutrino mang năng lượng động học lớn, gần như ngược chiều nhau, mỗi hạt có cỡ \(m_W/2\). Năng lượng nghỉ không hề biến mất; nó đã trở thành năng lượng của trạng thái cuối. Các kênh muon và tau, cũng như các kênh quark, cũng có thể xảy ra.

\begin{classroomqa}
\audiencequestion{Một \(W\) thực sống bao lâu, và có thể tính thời gian sống của nó từ cùng hằng số ghép đó không?
\lecturetimestamp{00:49:49}}
\lecturerresponse{Có. Ta tính biên độ phân rã cho từng kênh được phép, bình phương nó, lấy tích phân trên không gian pha của trạng thái cuối, rồi cộng các bề rộng riêng phần. Nghịch đảo của tổng bề rộng cho thời gian sống. Ước lượng số học được nói ra trong bản chép lời bị hỏng, nhưng ý định về mặt định tính là thời gian sống riêng là cực kỳ ngắn và một \(W\) đang chuyển động chịu giãn thời gian theo cách thông thường.\editorialnote{Dùng bề rộng đo được
\(\Gamma_W\simeq2.1\,{\rm GeV}\) cho
\(\tau_W=\hbar/\Gamma_W\simeq3.2\times10^{-25}\,{\rm s}\). Con số hiện đại này giải quyết đoạn trao đổi số liệu bị hỏng mà không gán nó cho bài giảng.}}
\end{classroomqa}

\begin{classroomqa}
\audiencequestion{Biên độ bình phương trở thành một tốc độ, và do đó thành một thời gian sống, như thế nào?
\lecturetimestamp{00:50:52}}
\lecturerresponse{Phần tử ma trận bình phương được kết hợp với mật độ các trạng thái cuối khả dĩ, các hàm delta bảo toàn và các hệ số chuẩn hóa để thu được xác suất trên một đơn vị thời gian, tức bề rộng phân rã \(\Gamma\). Xác suất sống sót là \(e^{-\Gamma t}\) trong các đơn vị với \(\hbar=1\), nên \(\tau=1/\Gamma\), hay \(\tau=\hbar/\Gamma\) khi khôi phục đơn vị. Vì thế, bề rộng yếu nhỏ đồng nghĩa với thời gian sống yếu dài.}
\end{classroomqa}

\begin{classroomqa}
\audiencequestion{Có thể gán một vận tốc cho hạt đang xuyên hầm qua rào không?
\lecturetimestamp{00:52:42}}
\lecturerresponse{Ta có thể lấy một khoảng cách suy ra chia cho một thời gian đi qua suy ra, nhưng đại lượng đó không nên được diễn giải là vận tốc của một hạt cổ điển hay của một tín hiệu. Việc cố đo một quỹ đạo bên trong rào làm thay đổi trạng thái và phá hủy thí nghiệm xuyên hầm ban đầu.}
\end{classroomqa}

\section{Quan sát một hạt ảo có nghĩa là gì?}

\subsection{Hạt ngoài và đường trong}

Hạt thực là các trạng thái tiệm cận đi vào và đi ra có thể được chuẩn bị hoặc phát hiện ở xa vùng tương tác. Hạt ảo là các đường trong một biểu diễn nhiễu loạn của biên độ. Chúng không phải là những vệt hạt ẩn bổ sung trong máy dò.

\begin{classroomqa}
\audiencequestion{Có định nghĩa rõ ràng nào về sự khác nhau giữa hạt thực và hạt ảo không?
\lecturetimestamp{00:55:05}}
\lecturerresponse{Hạt thực là một trạng thái ngoài ở trên vỏ khối đi vào hoặc đi ra khỏi vùng tương tác. Hạt ảo là một đường trong ở ngoài vỏ khối trong giản đồ Feynman. Đối tượng sau là một phần của khai triển toán học, không phải một vật thể được quan sát độc lập.}
\end{classroomqa}

Để phân giải một quá trình chỉ kéo dài \(\Delta t\), một đầu dò cần các tần số cỡ \(1/\Delta t\) và do đó năng lượng cỡ \(\hbar/\Delta t\). Một photon đủ sắc nét để phân giải một hạt alpha bên dưới một rào cao có thể cung cấp năng lượng cần thiết để nâng nó vào vùng rào. Tương tự, một đầu dò đủ năng lượng để phân giải \(W\) nội sống ngắn có thể tạo ra một \(W\) thực. Nỗ lực quan sát ấy biến quá trình ban đầu thành một thí nghiệm khác.

Bài giảng diễn đạt điều này như một thế lưỡng nan lặp đi lặp lại của cơ học lượng tử: sự can thiệp cần thiết để kiểm tra bức tranh trung gian bị cấm lại cung cấp chính năng lượng khiến cấu hình tưởng như bị cấm trở thành thực.

\subsection{Thăng giáng của proton}

Trong lý thuyết nhiễu loạn, một trạng thái proton chứa các biên độ cho những cấu hình như
\begin{equation}
  p\longleftrightarrow n+W^+,
\end{equation}
và một neutron tương tự cũng có
\begin{equation}
  n\longleftrightarrow p+W^-.
\end{equation}
Nếu không có tương tác nào khác xen vào, boson nội sẽ bị tái hấp thụ. Thay vào đó, một đầu dò cứng có thể để lại một \(W\) trong số các hạt đi ra. Ta có thể gọi hành động đó là đánh bật ra một thăng giáng hoặc là tạo ra một \(W\). Ngôn ngữ thì khác nhau; phát biểu có thể đo được là cùng một biên độ trạng thái cuối.

\begin{classroomqa}
\audiencequestion{Khi một đầu dò cứng tạo ra một \(W\) đi ra, nó đã đánh bật ra một hạt vốn đã ở đó hay tạo ra một hạt mới?
\lecturetimestamp{01:01:20}}
\lecturerresponse{Cả hai bức tranh đều có thể dùng để tổ chức toán học. Một mô tả sắp theo thứ tự thời gian nói rằng proton đã thăng giáng thành \(n+W^+\) và đầu dò lấy neutron đi trước khi xảy ra tái hấp thụ; một mô tả khác nói rằng va chạm đã tạo ra \(W\) ở trạng thái cuối. Điều bất biến là xác suất được tính cho trạng thái cuối quan sát được, chứ không phải một câu chuyện cổ điển duy nhất về cái gì đã tồn tại trước phép đo.}
\end{classroomqa}

Việc tăng năng lượng va chạm không phơi bày ra một bảng kê cố định các đối tượng vốn đã ẩn bên trong bia. Nó mở ra nhiều trạng thái cuối khả dĩ: một \(W\), vài boson, các jet quark, hadron và các hạt khác được cho phép bởi các số lượng tử bảo toàn. Ngay cả khi tác động đủ mạnh lên một trường sơ cấp, ta cũng có thể tạo ra proton và nhiều vật thể hợp phần khác. Thứ đi ra không nhất thiết đã hiện diện sẵn từ trước như một thành phần. Bài giảng đẩy ý này đến mức hài hước: nếu năng lượng và các số lượng tử cho phép thì ngay cả bàn, ghế, hay một nguyên tử hydro cũng không bị loại trừ về mặt logic như những sắp xếp trạng thái cuối cực kỳ khó xảy ra. Trò đùa ấy mang theo một lời cảnh báo nghiêm túc chống lại việc định nghĩa tính sơ cấp bằng một bảng kê ngây thơ các mảnh vỡ va chạm.

\subsection{Thăng giáng chân không và nhiệt độ}

Một thăng giáng gắn với proton không giống với một giản đồ chân không. Cái thứ nhất là một thành phần của trạng thái proton có tương tác; còn một thăng giáng chân không thì không có đường proton ngoài. Trong cả hai trường hợp, một đầu dò đủ năng lượng đều có thể tạo ra các hạt thực tương ứng với một trạng thái cuối được phép. Chẳng hạn, hai photon năng lượng cao có thể tạo ra một cặp electron--positron khi năng lượng trong hệ khối tâm vượt ngưỡng.

\begin{classroomqa}
\audiencequestion{Những phản ứng này có phụ thuộc vào nhiệt độ không?
\lecturetimestamp{01:05:17}}
\lecturerresponse{Biên độ của quá trình sơ cấp không đòi hỏi một bể nhiệt.
Nhiệt độ có ý nghĩa về mặt thống kê: một hệ ở nhiệt độ cao chứa nhiều
hạt có động năng lớn, nên các va chạm vượt ngưỡng trở nên
phổ biến. Việc nung nóng một hệ có chứa một hạt alpha liên kết cũng có thể cung cấp
đủ năng lượng cho sự thoát vượt rào, ngoài xuyên hầm lượng tử.
Một vài hạt giàu năng lượng cô lập tạo thành một chùm; nhiều hạt giàu năng lượng
được phân bố nhiệt tạo thành một bể nóng.}
\end{classroomqa}

\begin{classroomqa}
\audiencequestion{Hạt \(W\) được phát ra trong một quá trình thăng giáng
của proton có phải là một thăng giáng của chân không không?
\lecturetimestamp{01:06:56}}
\lecturerresponse{Không. \(W\) nội tại trong
\(p\leftrightarrow n+W^+\) là một thăng giáng của trạng thái proton tương tác.
Một quá trình ảo khép kín không có proton ngoài là một thăng giáng chân không.
Chúng tuân theo cùng các quy tắc của lý thuyết trường nhưng thuộc về những biên độ khác nhau.}
\end{classroomqa}

\begin{classroomqa}
\audiencequestion{Vì sao phản neutrino là cần thiết tại đỉnh leptonic của \(W\)?
\lecturetimestamp{01:08:51}}
\lecturerresponse{Nó mang các số lượng tử lepton được bảo toàn và là một phần
của đỉnh dòng tích điện cơ bản nối \(e\), \(\nu_e\) và \(W\).
Dấu gạch trên phân biệt phản neutrino electron đi ra với neutrino.}
\end{classroomqa}

\section{Có thể đẩy sự chuyển hóa năng lượng đến đâu?}

Phần tiếp theo của bài giảng đẩy quy tắc rằng năng lượng có thể tạo ra các hạt đến một
giới hạn cố ý phi lý. Một hạt cà phê có thể trở thành một chiếc bánh sô-cô-la không?
Không có quá trình thực tế nào kiểu đó, nhưng lý thuyết lượng tử tương đối tính
không áp đặt quy tắc rằng vật chất ở trạng thái cuối phải giống với vật ban đầu. Nó
đòi hỏi đủ năng lượng, các điện tích bảo toàn thích hợp, và một biên độ khác không.
Một kilôgam có năng lượng cỡ
\begin{equation}
  Mc^2\sim9\times10^{16}\ {\rm J},
\end{equation}
nên bất kỳ va chạm vi mô nào có thể lắp ráp nên một vật cụ thể ở thang kilôgam
đều sẽ có xác suất nhỏ đến mức không tưởng.

Khối lượng Planck,
\begin{equation}
  m_{\rm P}=\sqrt{\frac{\hbar c}{G}}
  \simeq2.2\times10^{-5}\ {\rm g},
\end{equation}
không phải là khối lượng lớn nhất có thể có. Nó là thang mà tại đó hấp dẫn trở nên
mạnh đối với một lượng tử được định xứ ở bước sóng Compton của chính nó. Một kilôgam chứa
xấp xỉ \(5\times10^7\) khối lượng Planck.

\begin{classroomqa}
\audiencequestion{Nếu các va chạm vượt quá thang Planck, các vật nặng tùy ý
như ghế hay hành tinh có thể xuất hiện không?
\lecturetimestamp{01:10:37}}
\lecturerresponse{Ở năng lượng khối tâm siêu-Planck tập trung trong một
vùng nhỏ, kết cục tự nhiên áp đảo là sự sụp đổ hấp dẫn:
một lỗ đen. Một vật vĩ mô có tổ chức rất cao sẽ là một
thăng giáng cực kỳ hiếm. Bài giảng viện dẫn phép đảo ngược thời gian: vì
một cái ghế có thể rơi vào lỗ đen, nên biên độ của quá trình ngược lại không hoàn toàn
bằng không. Xác suất của nó bị ức chế bởi một hệ số entropy khổng lồ, nên
không được nhầm lẫn giữa tính khả hữu và khả năng xảy ra trên thực tế.}
\end{classroomqa}

Sự bàn luận ngoài lề này củng cố lời cảnh báo trước đó. Lý thuyết trường lượng tử dự đoán
các biên độ giữa những trạng thái đầu và cuối được chỉ định. Một khẳng định
mang tính hình tượng rằng mọi sản phẩm theo nghĩa đen đều đã nằm sẵn trong bia trước khi va chạm
không phải là một phần của dự đoán đó.

\section{Boson, tương tác của chúng và thống kê của chúng}

\subsection{Photon có thể tán xạ}

Là boson không có nghĩa là không tương tác. Boson hay fermion quy định thống kê trao đổi; còn các đỉnh tương tác là những thực tế động lực học riêng biệt. Các chùm sáng cổ điển thường đi xuyên qua nhau vì QED không có đỉnh bốn photon ở mức cây và tán xạ photon--photon năng lượng thấp là cực kỳ nhỏ. Tuy vậy, hai photon có thể tương tác thông qua một vòng electron--positron ảo:
\begin{equation}
  \gamma+\gamma
  \longrightarrow(e^-e^+)_{\rm virtual}
  \longrightarrow\gamma+\gamma.
\end{equation}
Khi năng lượng khối tâm đủ cao, thay vào đó chúng có thể tạo ra một cặp thực,
\begin{equation}
  \gamma+\gamma\longrightarrow e^-+e^+.
\end{equation}
Các photon quang học có năng lượng cỡ vài electron volt, trong khi ngưỡng tạo cặp có bậc \(2m_e c^2\approx1.02\,{\rm MeV}\). Vì vậy, quá trình ảo bị ức chế mạnh ở các năng lượng quang học thông thường, đó là lý do các chùm sáng cắt nhau dường như không nhận ra nhau.

\begin{classroomqa}
\audiencequestion{Có phải các boson \(W\) ảo tương tác với các boson khác đơn giản chỉ vì chúng là boson không?
\lecturetimestamp{01:12:47}}
\lecturerresponse{Chúng có thể tương tác, nhưng không phải chỉ vì chúng tuân theo thống kê Bose. Các đỉnh chuẩn quyết định các tương tác động lực học. Còn thống kê quyết định cách các biên độ nhiều hạt ứng xử khi hoán đổi các hạt đồng nhất. Photon tương tác yếu thông qua các hạt ảo mang điện, và lý thuyết điện yếu cũng chứa các tự tương tác trực tiếp phi Abel của các boson \(W\).}
\end{classroomqa}

\subsection{Tăng cường Bose không phải là một lực}

Sự khác biệt giữa thống kê fermion và boson trở nên nổi bật khi nhiều hạt đồng nhất chiếm các trạng thái lượng tử lân cận. Boson thiên về việc nhiều hạt cùng chiếm một trạng thái; fermion thì loại trừ điều đó. Trong một va chạm năng lượng cao điển hình, các hạt bay vào những trạng thái động lượng tách biệt rõ ràng, nên hiệu ứng thống kê này có thể kém lộ rõ hơn so với trong laser hay chất ngưng tụ.

Về nguyên tắc, các boson \(W\) cũng có cùng sự tăng cường Bose. Bài giảng đùa rằng ta có thể tưởng tượng một hệ phát laser các boson \(W\). Khối lượng khổng lồ và thời gian sống cực ngắn của chúng khiến một thiết bị như vậy phi thực tế đến mức khó tin, nhưng thống kê của chúng quả thật là thống kê Bose.

\begin{classroomqa}
\audiencequestion{Proton hay neutron có màu gì?
\lecturetimestamp{01:18:41}}
\lecturerresponse{Một hadron vật lý là một singlet màu, thường được gọi là không màu hoặc trắng. Cách nói tắt là một baryon chứa một quark đỏ, một quark lục và một quark lam. Chính xác hơn, hàm sóng màu của nó là chồng chập phản đối xứng
\[
  |{\rm singlet}\rangle
  =\frac1{\sqrt6}\epsilon_{abc}|q^a q^b q^c\rangle,
\]
chứ không phải một sự gán theo kiểu cổ điển, trong đó mỗi quark được nêu tên mang một màu cố định vĩnh viễn.}
\end{classroomqa}

\begin{classroomqa}
\audiencequestion{Khóa học này có đề cập đến ngưng tụ Bose không?
\lecturetimestamp{01:19:25}}
\lecturerresponse{Có, thông qua hiện tượng Higgs. Một trường Higgs khác không khắp chân không tương tự, theo những phương diện quan trọng, với một ngưng tụ boson. Phép loại suy này sẽ được phát triển sau khi đã giới thiệu sự phá vỡ đối xứng tự phát.\editorialnote{Giá trị kỳ vọng chân không của Higgs không phải, theo nghĩa đen, là một khí thông thường gồm các hạt Higgs on-shell cùng chiếm một trạng thái. Sự ngưng tụ là một phép loại suy soi sáng cho một trường boson kết hợp và một chân không được chọn, chứ không phải một sự đồng nhất vi mô hoàn toàn.}}
\end{classroomqa}

\section{Boson yếu trung hòa còn thiếu}

Giờ hãy quay lại bảng nhóm. Nếu đối xứng yếu ngang là \(SU(2)\),
thì nó có
\begin{equation}
  \dim SU(2)=3
\end{equation}
các bộ sinh và do đó có ba trường chuẩn. Hai tổ hợp mang điện
\(W^+\) và \(W^-\)
chỉ tương ứng với hai trường. Cần có một trường thứ ba, trung hòa.
Nó liên hệ mật thiết với boson \(Z\) đã quan sát được, nhưng không phải đơn giản là
photon. Các tổ hợp trung hòa chính xác bao gồm cả trường thứ ba của \(SU(2)_L\)
lẫn trường \(U(1)_Y\). Sự trộn đó, cùng với định nghĩa đúng của điện tích,
được chủ ý để dành sang bài giảng tiếp theo.

\section{Phá vỡ đối xứng tự phát}

\subsection{Một lựa chọn đối xứng có thể ở trong thế lưỡng lự không ổn định}

Hãy hình dung một bàn ăn tròn mà ở mỗi chỗ ngồi có một dụng cụ ăn ở
bên trái và một ở bên phải. Trước khi có ai cử động, trái và phải là
tương đương. Một khi một vị khách với sang phải, mọi lựa chọn lân cận đều bị
ép buộc và cả bàn ổn định vào cấu hình mà mọi người đều chọn bên phải. Các
quy tắc vốn đối xứng, nhưng cấu hình được hiện thực hóa đã chọn một trong hai
khả năng. Câu chuyện chỉ là một phép loại suy, nhưng nó tách ra đúng tổ hợp
mà ta cần: đối xứng, nhiều kết cục tương đương, và một tính không ổn định cho
phép một nhiễu loạn nhỏ tùy ý chọn một kết cục trên toàn bộ hệ.

\subsection{Mạng đầu--đuôi}

Để có một mô hình gọn hơn, hãy đặt một biến hai trạng thái \(s_i=\pm1\) tại mỗi nút của
một mạng. Gọi \(+1\) là mặt ngửa và \(-1\) là mặt sấp. Giả sử các nút lân cận ưa
đồng hướng:
\begin{equation}
  H=-J\sum_{\langle ij\rangle}s_i s_j,
  \qquad J>0.
\end{equation}
Các láng giềng song song đóng góp năng lượng thấp hơn các láng giềng phản
song song. Vì thế có hai trạng thái cơ bản suy biến,
\begin{equation}
\begin{aligned}
  s_i&=+1\quad\text{for every }i,
  \\
  s_i&=-1\quad\text{for every }i.
\end{aligned}
\end{equation}
Hamiltonian bất biến dưới phép lật toàn cục \(s_i\to-s_i\), nhưng mỗi trạng thái
cơ bản được chọn đều không bất biến.

Bây giờ hãy cố định một spin ở biên rất xa thành mặt ngửa. Trong một hệ
thông thường có tương tác tầm ngắn, ta có thể kỳ vọng một bậc tự do ở xa như
vậy sẽ không ảnh hưởng gì ở đây. Tuy nhiên, trong pha có trật tự, điều kiện
biên đó chọn chân không toàn mặt ngửa. Một khi đã được chọn, một mặt
sấp đơn lẻ trong biển mặt ngửa sẽ tốn năng lượng vì nó không phù hợp với các
láng giềng của nó. Các kích thích cục bộ được xác định tương đối so với trạng
thái cơ bản đã chọn, mặc dù các phương trình nền tảng vẫn giữ đối xứng mặt
ngửa--mặt sấp.

\subsection{Phá vỡ tường minh và vách miền}

Phá vỡ tường minh sẽ thêm vào một thiên lệch cục bộ,
\begin{equation}
  H_{\rm explicit}
  =-J\sum_{\langle ij\rangle}s_i s_j-h\sum_i s_i,
  \qquad h\ne0.
\end{equation}
Khi đó, mặt ngửa và mặt sấp không có cùng năng lượng ngay cả tại một vị trí cô
lập, và chỉ còn một chân không được ưu tiên. Trong phá vỡ tự phát \(h=0\):
các quy tắc cục bộ là đối xứng và hai chân không vẫn suy biến, nhưng hệ chiếm
một trong số đó.

Thay vào đó, hãy áp đặt rằng biên ngoài cùng bên trái toàn là ngửa và biên ngoài cùng bên phải toàn là sấp. Cấu hình năng lượng thấp nhất phải nội suy giữa hai chân không đó. Mặt phân cách là một vách miền.

\begin{figure}[htbp]
\centering
\includegraphics[width=0.9\textwidth]{lecture_06_figure_08.png}
\caption{Mạng ngửa--sấp được chia bởi một vách miền nét đứt. Các điều kiện biên đối nghịch buộc hệ phải chứa một mặt phân cách giữa hai trạng thái có trật tự suy biến.}
\lectureframe{01:30:35}
\end{figure}

\begin{classroomqa}
\audiencequestion{Điều gì phân biệt sự phá vỡ đối xứng tự phát với sự phá vỡ đối xứng tường minh thông thường?
\lecturetimestamp{01:28:36}}
\lecturerresponse{Với phá vỡ tường minh, Hamiltonian cục bộ đã gán các năng lượng khác nhau cho ngửa và sấp, nên không tồn tại cặp trạng thái cơ bản liên hệ với nhau bởi đối xứng. Với phá vỡ tự phát, các phương trình vẫn giữ đối xứng, các chân không là suy biến, và điều kiện biên hoặc một nhiễu động nhỏ tùy ý sẽ chọn ra một trong số đó. Khả năng tồn tại một vách miền nội suy giữa các chân không phân biệt là dấu hiệu nhận biết của trường hợp sau.}
\end{classroomqa}

Mô hình này cố ý ở mức sơ khai, nhưng nó chứa logic cốt lõi sẽ được chuyển sang trường Higgs: một lý thuyết đối xứng, một cấu trúc chân không bất ổn hoặc suy biến, một trạng thái cơ bản được chọn ra, và các hệ quả vật lý do lựa chọn đó quyết định.

\section{Điểm dừng hiện tại của bài giảng}

Bảng fermion biến nhóm của Mô hình Chuẩn thành một tập hợp các hướng vật lý. Đối xứng màu \(SU(3)\) chạy theo chiều dọc qua đỏ, lục và lam; các phép biến đổi yếu mang điện chạy theo chiều ngang qua các bộ đôi quark và lepton; một đối xứng pha cung cấp thừa số \(U(1)\) ban đầu. Sự phát xạ boson chuẩn hiện thực hóa các phép biến đổi nhóm này thành các tương tác.

Sự yếu của phân rã beta khi đó được quy không phải cho \(g_w\) nhỏ một cách đặc biệt, mà cho hạt trung gian nặng. Trong không gian khoảng cách, khối lượng làm hàm truyền bị cắt cụt ở ngoài thang Compton \(\hbar/(m_Wc)\). Trong không gian xung lượng,
\begin{equation}
  \mathcal A_{\rm weak}
  \sim\frac{g_w^2}{k^2+m_W^2}
  \simeq\frac{g_w^2}{m_W^2}
\end{equation}
ở các xung lượng cỡ hạt nhân. Các thí nghiệm năng lượng thấp đời đầu đã đo tổ hợp này dưới dạng hằng số Fermi.

Phần thảo luận lớp học mở rộng làm rõ công thức này có nghĩa gì và không có nghĩa gì. Các đường \(W\) nội bộ là ảo và off shell; năng lượng--xung lượng vẫn được bảo toàn. Một đầu dò đủ mạnh có thể tạo ra một \(W\) thực, nhưng khi đó nó đã tạo ra một quá trình với trạng thái cuối khác. Các sản phẩm va chạm không phải là một danh sách theo nghĩa đen của các thành phần tồn tại sẵn. Nhiệt độ làm thay đổi số lượng các hạt tới năng lượng cao, photon có thể tương tác thông qua các hạt ảo mang điện, và thống kê Bose là điều khác hẳn với một lực động lực học.

Cuối cùng, \(SU(2)\) đòi hỏi một trường chuẩn trung hòa còn thiếu liên hệ với \(Z\), và mạng ngửa--sấp cung cấp mô hình đầu tiên về sự phá vỡ đối xứng tự phát. Bước tiếp theo là kết hợp hai mạch còn dang dở đó: sự trộn điện yếu và cơ chế Higgs.
